%! Author = zhiweizhang
%! Date = 10.06.25

\chapter{Discussion and Analysis}\label{chapter:discussion-and-analysis}

This chapter reflects critically on the performance of the evaluated models, connecting the empirical results to the
research objectives outlined in Chapter~\ref{chapter:introduction}. The primary aim of this thesis was to develop
and compare image classification models capable of accurately recognizing and categorizing Nazi symbols from varied
visual inputs. Through detailed analysis of both binary and multi-class classification results, this chapter
interprets key trends, identifies strengths and limitations of different model families, and discusses the practical
and ethical implications of deploying such systems.

\section{Analysis of Binary Classification Performance}

The results presented in ~\autoref{tab:binary-results} reveal a wide range of performances across different models.
Notably, classical \ac{ML} classifiers leveraging OpenCLIP embeddings consistently ranked among the top
performers. The best overall F1-score (0.939) was achieved by the \textbf{OpenCLIP + ~\ac{SVC}} model, which also
demonstrated high precision (0.948) and recall (0.929). Similarly, the \textbf{OpenCLIP + Logistic Regression}
model achieved the highest ~\ac{AUC} score (1.000) and highest \ac{AP} (0.964), along with a strong recall of 0.882
and an F1-score of 0.902, indicating a solid ability to correctly classify Nazi symbols while minimizing false
negatives.

Among \ac{DL} models, both \textbf{\ac{ViT}} and \textbf{\ac{YOLO}} performed competitively. \ac{ViT} achieved an
F1-score of 0.901, the highest among all standalone neural models, while \ac{YOLO} achieved an F1-score of 0.857 with
a notably high recall of 0.932, highlighting its capacity to correctly classify positive symbol instances.
\ac{ResNet} also performed well, with an F1-score of 0.867, slightly outperforming EfficientNet (0.767).

Anomalously, the \textbf{Moondream} model achieved an exceptionally high recall of 0.995 but suffered from very low
precision (0.153), resulting in a low F1-score of just 0.265. This suggests a strong tendency to overpredict the
Nazi class, leading to a high number of false positives. Similarly, the \textbf{OpenCLIP zero-shot} model performed
poorly, with a low F1-score of 0.025, highlighting the limitations of zero-shot classification without additional
fine-tuning or classifier integration.

Several classical ensemble models also showed robust performance when paired with OpenCLIP embeddings. For instance,
\textbf{OpenCLIP + VotingClassifier} reached an F1-score of 0.908 and precision of 0.945, while
\textbf{OpenCLIP + AdaBoostClassifier} and \textbf{OpenCLIP + GradientBoostingClassifier} achieved F1-scores of
0.842 and 0.838, respectively. However, the \textbf{OpenCLIP + RandomForestClassifier} produced an unusually high
precision (0.983) but very low recall (0.097), resulting in a poor F1-score (0.177), indicating that it was highly
conservative in predicting the Nazi class.

Overall, models that combine OpenCLIP embeddings with classical classifiers—especially \ac{SVC}, KNeighborsClassifier, and
VotingClassifier—outperformed all other approaches in this binary classification task. Nonetheless, end-to-end deep
learning models like \ac{ViT} and \ac{YOLO} remain competitive and demonstrate strong generalization capabilities,
particularly when trained on high-resolution inputs and sufficient labeled data.

\subsection{Binary Classification Confusion Matrix Analysis}

~\autoref{fig:confusion-matrix-binary} presents the confusion matrices for the binary classification task, where the
goal is to distinguish between \textit{Nazi} and \textit{Non-Nazi} images.

The OpenCLIP encoder combined with \ac{SVC} achieves the best overall balance, producing the
highest number of true positives (194) and true negatives (14,802), alongside the lowest false positives (11) and
false negatives (11). This indicates a strong capability to correctly classify Nazi symbols while minimizing
misclassifications.

\ac{YOLO} shows excellent sensitivity, with the second lowest false negative count (14), meaning it rarely misses Nazi
instances, but it has a relatively higher false positive rate (50), which might lead to more false alarms.

\ac{ResNet} also performs well, showing a favorable trade-off between false positives (34) and false negatives (22),
maintaining a good level of accuracy and reliability.

\ac{ViT} provides balanced performance, with moderate false positives (28) and false negatives (18), showing
robustness in both detecting Nazi symbols and correctly rejecting non-Nazi images.

EfficientNet, while demonstrating strong true negative results, has the highest false negatives (49), indicating
that it misses more Nazi cases compared to the other models, which might limit its applicability when recall is critical.

Overall, these confusion matrices confirm that the OpenCLIP + \ac{SVC} approach is the most effective for binary
classification in this context, offering a strong balance between precision and recall essential for sensitive
symbol classification.

\subsection{Binary Classification \ac{ROC} and \ac{PR} Curve Analysis}

The \ac{ROC} curve results indicate that most models achieved excellent discriminative ability in the binary
classification task, with several models reaching near-perfect \ac{AUC} values close to 1.0. Notably, OpenCLIP combined
with \ac{LR}, KNeighborsClassifier, and Random Forest classifiers, along with \ac{ViT}, \ac{YOLO},
and \ac{ResNet}, all demonstrated near-identical, near-optimal \ac{ROC} performance (\ac{AUC} $\approx$ 1.000),
suggesting their strong capacity to separate Nazi and Non-Nazi classes.

Ensemble methods such as AdaBoost and Gradient Boosting classifiers paired with OpenCLIP embeddings also performed
robustly (\ac{AUC} $\approx$ 0.998), while EfficientNet’s slightly lower \ac{AUC} of 0.993 still represents a high-quality
classification boundary. Bagging and Decision Tree classifiers showed comparatively lower \ac{AUC}s (0.955 and 0.933),
indicating reduced robustness and potential susceptibility to misclassification in borderline cases.

The \ac{AP} values from the \ac{PR} curves further underscore these findings, revealing that
OpenCLIP + \ac{LR} and \ac{ViT} provide the best balance between precision and recall, critical in scenarios where false
positives and false negatives carry significant consequences. The lower \ac{AP} of the OpenCLIP + Decision Tree
classifier (0.625) suggests that, despite reasonable \ac{AUC} scores, this model struggles to maintain precision at
higher recall levels, reflecting less reliable positive class predictions.

These combined results imply that embedding-based methods (OpenCLIP) paired with strong classifiers (\ac{LR},
KNeighborsClassifier) or transformer architectures (\ac{ViT}) are most effective for the task. In contrast, simpler or single-model
approaches such as Decision Trees show limitations in practical use despite acceptable \ac{ROC} metrics.

Only models that support probabilistic outputs are included in this analysis. Classifiers such as \texttt{\ac{SVC}},
\texttt{SGDClassifier}, and the \texttt{VotingClassifier}, as well as models like \texttt{Moondream} and
\texttt{OpenCLIP}, are omitted from \ac{ROC} and \ac{PR} curve evaluations, as discussed in
~\autoref{sec:evaluation-metrics}.

\subsection{Binary Classification with Reduced Nazi Symbol Classes}

To further examine model robustness and adaptability, a binary classification experiment was conducted. This
experiment was restricted to a reduced subset of Nazi symbols (\textit{siegrune}, \textit{swastika}, and
\textit{\ac{SS} skull}) against the full set of non-Nazi images. The goal was to evaluate model performance when the
semantic diversity of the positive class is reduced.

Despite the narrower scope, overall trends remained consistent with the full binary task (as showed in
~\autoref{tab:binary-reduced-symbols-result}). Transformer-based
architectures continued to perform strongly. \ac{ViT} achieved excellent results with an F1-score of \textbf{0.948},
indicating its strong ability to generalize even with reduced symbol variety. Among OpenCLIP-based hybrid models,
\textit{OpenCLIP + \ac{SVC}} achieved the best overall performance with an accuracy of \textbf{0.998},
precision of \textbf{0.973}, recall of \textbf{0.900}, and F1-score of \textbf{0.935}, while
\textit{OpenCLIP + KNeighborsClassifier} attained the highest recall at \textbf{0.989}.

Interestingly, \ac{YOLO} underperformed in this setup (F1-score \textbf{0.661}) compared to its results on the full
binary task, potentially due to its reliance on spatial object features that may be less effective when class
variety is reduced. Conversely, models like \textbf{Moondream} and \textbf{OpenCLIP} hybrids maintained high performance, indicating
that text-image aligned embeddings may still effectively capture distinguishing features in constrained settings.

As this is a focused experiment with a reduced symbol set, the evaluation was limited to the most informative and widely
supported metrics—accuracy, precision, recall, and F1-score. Metrics such as \ac{AUC} and \ac{AP} are omitted for
simplicity and consistency, especially given that not all models provide calibrated probability outputs.

This experiment highlights that reducing the complexity of the positive class does not necessarily simplify the task
uniformly for all models. While it benefits embedding-based methods and transformers, object detection models like
\ac{YOLO} may be affected due to loss of visual diversity.

\section{Analysis of Multi-Class Classification Performance}

~\autoref{tab:multi-class-results} summarizes the performance of all evaluated models on the multi-class
classification task. The \textbf{\ac{YOLO}} model achieved the best performance across all macro-averaged metrics,
with an accuracy of 0.884, precision of 0.865, recall of 0.818, and F1-score of 0.833. This superior performance
likely stems from \ac{YOLO}’s capability to extract fine-grained spatial features in high-resolution images, combined
with its efficient end-to-end object-centric architecture.

The \textbf{\ac{ResNet}} and \textbf{\ac{ViT}} models also demonstrated strong performance, with F1-scores of 0.810 and
0.790, respectively. \ac{ResNet} maintained a slightly higher accuracy (0.879) than \ac{ViT} (0.876), suggesting
that deep convolutional features still offer competitive advantages. Meanwhile, \ac{ViT}’s transformer-based design
effectively captured semantic dependencies, reinforcing its value for tasks involving subtle symbol variations.

In contrast, \textbf{EfficientNet} underperformed in this task, with an F1-score of 0.709 and a macro recall of
0.705, indicating difficulty in generalizing across all classes. The \textbf{Moondream} model also showed weak
results (F1-score of 0.355), despite a moderately high recall (0.601), suggesting overprediction of certain classes.

The \textbf{OpenCLIP zero-shot} model performed poorly overall, with an accuracy of 0.274 and F1-score of 0.265.
These results reinforce that zero-shot approaches struggle when applied to highly specialized tasks like Nazi symbol
recognition without task-specific adaptation or supervised training.

When OpenCLIP embeddings were combined with traditional classifiers, performance varied widely.
\textbf{OpenCLIP + KNeighborsClassifier} (F1-score: 0.775) and \textbf{OpenCLIP + \ac{SVC}} (F1-score: 0.784) were
the most effective, indicating that neighborhood-based and kernel-based methods are capable of partially exploiting
OpenCLIP’s semantic features. Similarly, \textbf{OpenCLIP + SGDClassifier} and
\textbf{OpenCLIP + GradientBoostingClassifier} achieved solid F1-scores of 0.718 and 0.703, respectively.

Conversely, ensemble-based models such as \textbf{RandomForest}, \textbf{Bagging}, and
\textbf{DecisionTreeClassifier} performed poorly, with F1-scores around or below 0.200. These results suggest that
these models may not align well with the high-dimensional structure of OpenCLIP embeddings, potentially due to
overfitting or limited ability to model fine-grained inter-class distinctions.

Overall, end-to-end models such as \ac{YOLO}, \ac{ResNet}, and \ac{ViT} clearly outperform both zero-shot approaches and
OpenCLIP embedding pipelines in the multi-class setting. These findings emphasize the importance of strong visual
feature extraction and task-aligned architectures when handling classification tasks involving subtle visual
differences across many categories.

\subsection{Comparison with Binary Classification Performance}

Comparing the results from the binary (~\autoref{tab:binary-results}) and multi-class
(~\autoref{tab:multi-class-results}) classification tasks reveals consistent trends in model effectiveness, but also
highlights some key differences driven by task complexity.

In both settings, \textbf{\ac{YOLO}} and \textbf{\ac{ViT}} emerged as top-performing \ac{DL} models. However, their
relative strengths became more pronounced in the binary task. For instance, \ac{ViT} achieved an F1-score of 0.901 and
\ac{YOLO} reached 0.857 in the binary case, whereas their multi-class F1-scores dropped to 0.790 and 0.833,
respectively. This drop reflects the increased difficulty of distinguishing between a larger number of visually
similar symbols.

The benefit of combining \textbf{OpenCLIP embeddings with classical classifiers} was more significant in the binary
setting. In that task, models like \textbf{OpenCLIP + \ac{SVC}} and \textbf{OpenCLIP + KNeighborsClassifier} achieved
F1-scores of 0.939 and 0.915 respectively, outperforming nearly all other approaches. However, in the multi-class
scenario, while still strong (F1-scores of 0.784 and 0.775 respectively), their advantage diminished, likely due to
increased inter-class ambiguity and the limitations of traditional classifiers in modeling fine-grained boundaries
across many classes.

Interestingly, some ensemble models (e.g., GradientBoosting, VotingClassifier) showed reasonable binary
classification performance (F1-scores around 0.840–0.910), but generally degraded in the multi-class case (F1-scores
below 0.710), possibly due to overfitting or insufficient class-separation power.

The \textbf{Moondream} and \textbf{OpenCLIP zero-shot} models performed poorly in both tasks, with Moondream showing
a pathological tendency to predict the positive class in binary classification (perfect recall but extremely low
precision), and generally low macro F1-scores in multi-class classification. These outcomes highlight the risks of
deploying zero-shot or generative models in specialized symbol recognition tasks without task-specific training.

Overall, binary classification appears more forgiving, allowing classical classifiers with pretrained embeddings to
compete closely with end-to-end \ac{DL} models. In contrast, the multi-class task places a heavier burden on
precise visual discrimination and inter-class separability, where deep models like \ac{YOLO} and \ac{ViT} retain a clear
edge. This underscores the importance of matching model complexity and architecture with task granularity and class
diversity.

\subsection{Multi-Class Classification Confusion Matrix Analysis}

The confusion matrices for the multi-class classification task reveal distinct patterns of model performance across
classes, reflecting varying degrees of class separation and misclassification.

\textbf{\ac{ResNet}} (see ~\autoref{fig:resnet-confusion-matrix-multiple}) shows strong true positive rates for
high-frequency classes such as \textit{Swastika} (808 correct predictions) and \textit{Siegrune} (247 correct), but
also notable confusion between visually similar classes. For example, \textit{SS Skull} is frequently misclassified as
\textit{Siegrune} (36 instances), while \textit{Siegrune} mislabels \textit{SS Skull} less often (5 instances). This
asymmetric confusion highlights the difficulty \ac{ResNet} faces in discriminating fine-grained symbol variations.
Low-frequency classes like \textit{Nazi Salute} perform poorly, with very few correct predictions.

\textbf{\ac{YOLO}} (see ~\autoref{fig:yolo-confusion-matrix-multiple}) generally maintains competitive accuracy for
dominant classes, correctly identifying 834 \textit{Swastika} instances and 222 \textit{Siegrune} samples. Confusion
between \textit{SS Skull} and \textit{Siegrune} remains substantial but somewhat more balanced compared to
\ac{ResNet}, with 25 \textit{SS Skull} instances misclassified as \textit{Siegrune} and 19 misclassifications in the
opposite direction. This indicates improved inter-class separation but ongoing challenges for these similar classes.

\textbf{EfficientNet} (see ~\autoref{fig:efficient-confusion-matrix-multiple}) exhibits a somewhat different confusion
profile, with lower misclassification of \textit{SS Skull} as \textit{Siegrune} (15 instances), but higher
misclassification of \textit{Siegrune} as \textit{SS Skull} (40 instances), showing an inverse asymmetry compared to
\ac{ResNet}. The true positive count for \textit{Siegrune} is comparatively lower (174), suggesting less confident
classification in this region of the feature space. Low-frequency classes continue to be problematic, with
\textit{Nazi Salute} achieving only a single correct prediction.

\textbf{\ac{ViT}} (see ~\autoref{fig:vit-confusion-matrix-multiple}) marginally improves performance for certain classes, such as
\textit{SS Skull} and \textit{Hitler}, with 188 and 175 correct predictions respectively. Confusion between
\textit{SS Skull} and \textit{Siegrune} remains, with 19 \textit{SS Skull} samples misclassified as \textit{Siegrune} and 32 in the
reverse direction, indicating persistent challenges in distinguishing subtle differences despite the transformer
architecture’s strengths.

\textbf{OpenCLIP Encoder combined with \ac{SVC}} (see ~\autoref{fig:openclip-encoder-confusion-matrix-multiple})
achieves a balanced confusion pattern. It correctly identifies 232 \textit{Siegrune} and 167 \textit{SS Skull}
instances. The confusion between \textit{SS Skull} and \textit{Siegrune} is distributed moderately (25 and 17
misclassifications respectively), showing a more symmetric error pattern compared to other models. This suggests
better inter-class separation in this challenging pair, although some misclassification persists. Additionally, the
model excels at maintaining tight clusters around true labels for high-frequency classes such as
\textit{Swastika} (823 correct) but struggles with low-frequency classes like \textit{Nazi Salute}, which has zero
correct predictions.

Overall, all models face difficulties with fine-grained classes like \textit{SS Skull} and \textit{Siegrune}, but
OpenCLIP+\ac{SVC} demonstrates a more balanced confusion profile. The results highlight the ongoing need to improve
discrimination for visually similar symbols, especially for classes with limited training data.

\section{Analysis of Underperforming Models}

While several models demonstrated strong classification performance, it is equally important to investigate those
that underperformed to understand their limitations and potential failure modes.

The \textbf{Moondream} model consistently exhibited pathological behavior, particularly in the binary classification
setting. Its extremely high recall (0.995) was accompanied by very low precision (0.153), leading to an overall
F1-score of just 0.265. This indicates that Moondream predicted nearly every image as a Nazi symbol, resulting in
excessive false positives. This failure mode may stem from its generative nature or misalignment between the model’s
text-conditioned outputs and the visual symbol classification task. Its weak performance in multi-class
classification (F1-score: 0.355) further confirms its unsuitability for tasks requiring high visual specificity
without fine-tuning.

Similarly, the \textbf{OpenCLIP zero-shot} model underperformed in both binary and multi-class tasks. In the binary
task, it produced an F1-score of just 0.025, suggesting that its generic semantic alignment is insufficient for
accurate Nazi symbol classification without supervised adaptation. In the multi-class task, OpenCLIP zero-shot achieved
only 0.265 F1-score and an accuracy of 0.274, further highlighting its limitations in recognizing fine-grained,
domain-specific symbol categories. These results emphasize the need for careful domain adaptation when deploying
foundation models for specialized tasks.

Some \textbf{classical ensemble models}, such as \textbf{OpenCLIP + RandomForestClassifier}, also performed
unexpectedly poorly despite high precision (0.983) in binary classification. The extremely low recall (0.097)
resulted in a poor F1-score (0.177), indicating an overly conservative classification threshold. In multi-class
classification, these models (RandomForest, Bagging, Decision Tree) also failed to produce competitive results (
F1-scores below 0.200), likely due to their limited capacity to model complex, high-dimensional OpenCLIP features or
handle subtle inter-class differences without overfitting.

Overall, these underperformances highlight the importance of model-task alignment. Generative or zero-shot models
may offer generality, but without fine-tuning, they can fail catastrophically on specific classification tasks.
Likewise, conventional tree-based models may struggle to effectively utilize dense, semantic embeddings like those
from OpenCLIP, especially in tasks with fine-grained categories.

These failures provide a contrast to the following high-performing approaches, which better align with the
classification objectives.

\section{Analysis of High-Performing Models}

Across both binary and multi-class classification tasks, several models demonstrated exceptional performance,
providing insight into the characteristics of effective architectures for Nazi symbol recognition.

The highest-performing binary classification model was \textbf{OpenCLIP + \ac{SVC}}, which achieved an F1-score of
0.939, precision of 0.948, and recall of 0.929. This model’s strength lies in its ability to exploit OpenCLIP’s
semantically-rich embeddings using a powerful kernel-based decision boundary. Its confusion matrix confirms minimal
false positives and negatives, demonstrating high reliability in correctly identifying both the presence and absence
of Nazi symbols. Likewise, \textbf{OpenCLIP + Logistic Regression} achieved the best AUC (1.000) and AP (0.964),
showing strong probabilistic calibration and balanced performance across decision thresholds.

In the multi-class classification task, \textbf{\ac{YOLO}} emerged as the top-performing model (F1-score: 0.833),
outperforming all others in terms of macro-averaged accuracy and recall. YOLO’s object detection backbone and
ability to preserve high-resolution spatial details proved highly effective in distinguishing subtle visual features
across symbol classes. Its relatively balanced confusion matrix also shows good per-class discrimination,
particularly on high-frequency categories like \textit{Swastika} and \textit{Siegrune}.

\textbf{\ac{ResNet}} and \textbf{\ac{ViT}} consistently performed well in both tasks. \ac{ResNet}’s deep
convolutional structure allowed it to learn robust visual features, resulting in an F1-score of 0.867 (binary) and
0.810 (multi-class). \ac{ViT} leveraged self-attention to capture semantic relationships, achieving an F1-score of
0.901 (binary) and 0.790 (multi-class). Their success across tasks demonstrates the generalizability of end-to-end
\ac{DL} models when properly trained with domain-specific labeled data.

In embedding-based approaches, \textbf{OpenCLIP + KNeighborsClassifier} also stood out. It delivered strong binary
performance (F1-score: 0.915) and competitive multi-class results (F1-score: 0.775), suggesting that distance-based
methods are particularly effective when operating on semantically meaningful representations.

These findings confirm that both end-to-end deep models and hybrid embedding-based approaches can be highly
effective when appropriately matched to the task. Kernel- and attention-based architectures, along with models that
retain spatial resolution or exploit global semantics, perform especially well on complex symbol classification
challenges.

Beyond individual model metrics, broader themes emerge across tasks and models, offering insight into effective
strategies for symbol classification.

\section{Key Insights and Thematic Patterns}
\label{subsec:key-insights}

Across both binary and multi-class classification tasks, several key insights emerged:

\begin{itemize}
    \item Vision-language embeddings from OpenCLIP, when combined with traditional machine learning classifiers,
    consistently delivered strong performance in both binary and multi-class tasks. These hybrid models showed
    robustness to visual variations and generalized well to stylized or contextually embedded symbols —
    outperforming solo deep learning models in many cases. However, solo OpenCLIP or Moondream models performed
    poorly, suggesting that raw multi-modal outputs are not directly effective for this task without a downstream
    classifier.
    \item The \ac{ViT} achieved consistently strong performance across both tasks. Its global
    attention mechanism allowed it to model subtle differences between symbols and perform well even in visually
    cluttered or stylized inputs, outperforming CNNs in several scenarios.
    \item YOLO, though traditionally designed for object detection, was repurposed and evaluated as a classifier
    in this research. It achieved moderate performance in both binary and multi-class classification tasks. Its
    performance on the reduced binary classification dataset was reasonable, though it did not match the
    top-performing models like \ac{ViT} or OpenCLIP with classical ML classifiers. Its relatively
    lightweight architecture and real-time inference potential, however, suggest it could still be useful in
    deployment scenarios where speed is critical and classification accuracy is acceptable.
    \item Classic CNNs like ResNet delivered strong baseline performance, particularly for clean and
    high-resolution inputs, but struggled with rarer symbols and visually noisy data. Similarly, EfficientNet was
    efficient but less effective for symbols with subtle intra-class variation (e.g., SS bolts vs. Wolfsangel).
    \item A key theme was the challenge of generalizing to real-world and historical contexts — including grainy
    images, scans, and stylized iconography. Models trained on curated datasets showed limitations here, suggesting
    a need for domain adaptation or contrastive pretraining to improve transferability.
    \item Finally, class imbalance significantly impacted performance, especially in the multi-class task. Rare
    symbols suffered from low recall and inconsistent predictions, highlighting the importance of balanced sampling,
    targeted augmentation, or few-shot learning approaches in future work.
\end{itemize}

Together, these insights highlight the trade-offs between model complexity, interpretability, and task
suitability—informing how future systems can balance accuracy with practicality and ethical deployment needs.

\section{Implications for Real-World Deployment}
\label{subsec:deployment-implications}

The findings suggest that hybrid approaches combining deep feature extraction with classical \ac{ML}
classifiers offer the most promise for deployment in real-world systems tasked with binary classification of Nazi
symbols in images. In particular, OpenCLIP combined with \ac{SVC} consistently achieve the highest F1-score in the
binary classification task. indicating strong performance in balancing precision and recall for accurately
classifying rare and sensitive Nazi symbols. These models also offer a good balance of performance, model
transparency, and compute efficiency—important factors for safety-critical applications such as content moderation
or archivalscreening.

However, real-world deployment would require addressing several practical issues. These include handling ambiguous
or degraded imagery, unseen symbol variants, and high false-positive costs in sensitive legal or moderation
contexts. The results suggest that a human-in-the-loop system may be appropriate—one where automated classification
is followed by expert verification for uncertain cases.

Moreover, compliance and interpretability matter: using simpler classifiers on top of OpenCLIP embeddings offers a
degree of explainability that deep end-to-end models may lack, aiding trust and regulatory acceptance.

These considerations directly reflect the initial research goals of building practical, safe, and transparent
classifiers for real-world symbol classification applications.

\section{Limitations and Risks}
\label{subsec:limitations}

While the system showed encouraging results, several limitations remain. Most notably, performance on minority
classes was inconsistent, with models tending to overfit to frequent, visually dominant symbols. This raises
fairness and completeness concerns in applications that require comprehensive classification of all symbol types.

Another risk lies in dataset bias: most training data was synthetically curated or sourced from known repositories,
which may not capture the full spectrum of symbol usage in obscure or modern reinterpretations. This could lead to
symbols being misclassified or overlooked in real-world use.

Finally, the model does not currently account for contextual cues—e.g., the difference between educational versus
propagandistic usage. While this was beyond the scope of the project, it remains a critical consideration for
responsible deployment. Over-reliance on image classification alone, without contextual understanding, could lead to
false accusations or misinterpretations.
